
\documentclass[11pt]{article}
\usepackage{series}
\usepackage{booktabs}
\usepackage{array}

\renewcommand\SeriesShort{Delta Projection}
\renewcommand\PartRoman{XI}
\renewcommand\PartArabic{11}
\renewcommand\PartsInSeries{XI}

\title{MTT-Corrected Propagators and UV Behaviour\\
\large Canonical Coherent Filtering from Fixed-Point Data}
\author{Peter Nero}
\date{April 2026}

\begin{document}
\maketitle

\begin{abstract}
The preceding delta-projection papers established a structural dictionary: many Dirac delta functions in physics are sharp limits of finite kernels, filters, shells, windows, or survivor-basin effects. The remaining execution-level problem is to determine which finite kernel is selected by Modal Triplet Theory (MTT), rather than introducing an arbitrary regulator. Building on the canonical coherent-kernel construction from fixed-point data, this paper studies the first concrete dynamical application: propagators. In a flat coherent chart, with positive quadratic operator \(L_m=-\Delta+m^2\) and MTT proper-time scale \(\tau>0\), the canonical corrected Euclidean propagator is
\[
        \Delta_\tau(k)=\frac{e^{-\tau |k|^2}}{|k|^2+m^2}.
\]
More generally, with a coherent spectral window \(B(A)=P\chi(A)e^{-\tau A/2}\chi(A)P\), the corrected covariance is
\[
        \Delta_{\mathrm{MTT}} = B(A)L_m^{-1}B(A)^\ast .
\]
We prove that \(\Delta_\tau\) recovers the ordinary propagator distributionally as \(\tau\downarrow0\), has finite coincident value in every finite dimension for \(\tau>0\), and yields Gaussian UV domination. For scalar polynomial perturbation theory in a fixed coherent sector, every Euclidean Feynman integral whose internal lines carry the canonical Gaussian filter is absolutely UV finite. The leading low-momentum correction is controlled by \(\tau |k|^2\), giving an effective coherence scale \(\Lambda_{\mathrm{eff}}\sim\tau^{-1/2}\). These results do not yet fix the numerical value of \(\tau\) in a physical sector; rather, they show that once fixed-point operator data determine \(\tau\), the propagator correction is no longer arbitrary.
\end{abstract}

\tableofcontents

\section{Purpose and claim discipline}

The earlier papers in this sequence established forward singular-limit statements of the form
\[
        K_\epsilon \longrightarrow \delta .
\]
That is mathematically useful but not sufficient for prediction. The predictive direction is
\[
        \text{MTT fixed-point data}
        \Longrightarrow
        K_{\mathrm{MTT}}
        \Longrightarrow
        \text{finite correction}.
\]
The companion paper on canonical coherent kernels supplied the first step: fixed-point operator data select a finite kernel through functional calculus. The present paper applies that kernel to propagators.

The paper proves a controlled statement in a flat Euclidean coherent chart. It does not claim to complete the full interacting Standard Model, quantum gravity, or nonperturbative MTT. It proves that the canonical MTT filter gives a specific UV-softened propagator and that this propagator has precise convergence, finiteness, and perturbative domination properties.

\begin{center}
\begin{tabular}{p{0.28\linewidth}p{0.62\linewidth}}
\toprule
Status & Claim \\
\midrule
Proved here & Gaussian-proper-time filtered propagators are finite at coincidence and recover the local propagator as \(\tau\downarrow0\). \\
Proved here & Euclidean scalar Feynman integrals with filtered internal lines are absolutely UV finite. \\
Interpreted here & \(\tau\) is the downstream proper-time/coherence scale inherited from MTT fixed-point data. \\
Not proved here & The numerical value of \(\tau\) in a real physical sector. \\
Not proved here & Full Lorentzian unitarity, gauge invariance, or Standard Model phenomenology. \\
\bottomrule
\end{tabular}
\end{center}


\section{Novelty and non-novelty}

The formula
\[
        \Delta_\tau(k)=\frac{e^{-\tau |k|^2}}{|k|^2+m^2}
\]
is not, by itself, a new regularization formula. Exponential and heat-kernel regulators are standard tools in mathematical physics and quantum field theory. The novelty claimed here is therefore not the bare functional form.

The novelty claim is the following selection statement:
\[
        \boxed{
        \text{the filter is not chosen to regulate a divergent integral; it is inherited from fixed-point data.}
        }
\]
In an MTT fixed-point regime, the data \((A,P,\chi,\tau)\) have independent meaning before any loop integral is evaluated:
\begin{itemize}
\item \(A\) is the linearized damping/geometric operator around the coherent fixed point;
\item \(P\) is the coherent-sector projector;
\item \(\chi(A)\) is the admissible spectral acceptance window associated with the gap and truncation regime;
\item \(\tau\) is the proper-time/coherence support scale of the evolve--project cycle.
\end{itemize}
Thus the finite propagator is not an arbitrary Gaussian regulator. It is the downstream covariance selected by the fixed-point sector.

Equivalently, the paper distinguishes three statements:
\begin{center}
\begin{tabular}{p{0.32\linewidth}p{0.55\linewidth}}
\toprule
Statement & Status \\
\midrule
The factor \(e^{-\tau k^2}\) is a familiar heat-kernel filter & not novel \\
Using such a factor as a regulator & not novel by itself \\
Deriving \((A,P,\chi,\tau)\) from MTT fixed-point data and then obtaining the filter & the execution-level novelty target \\
\bottomrule
\end{tabular}
\end{center}

This paper proves the third statement in the model fixed-point setting where the operator data are specified. It does not yet claim that the numerical value of \(\tau\) has been derived for the Standard Model or for quantum gravity.


\section{Canonical fixed-point input}

Let \(X\) denote a flat coherent chart, either \(\mathbb R^d\) or a torus \(T^d\). Let
\[
        A=-\Delta
\]
be the nonnegative coherent-chart Laplacian. Let \(P\) denote the coherent-sector projector inherited from fixed-point data, and let \(\chi(A)\) be an admissible spectral window. The canonical coherent filter is
\[
        B_{\tau,\chi}
        :=
        P\chi(A)e^{-\tau A/2}\chi(A)P ,
        \qquad \tau>0.
\]
The corresponding kernel is
\[
        K_{\tau,\chi}(x,y)
        =
        \langle x|B_{\tau,\chi}B_{\tau,\chi}^{\ast}|y\rangle .
\]

In the simplest full flat chart, \(P=I\) and \(\chi=I\), so
\[
        B_\tau=e^{-\tau A/2},
        \qquad
        K_\tau=e^{-\tau A}.
\]
In momentum variables,
\[
        \widehat{K_\tau}(k)=e^{-\tau |k|^2}.
\]
Thus \(K_\tau\) is the heat-kernel replacement of the Dirac delta:
\[
        K_\tau(x-y)
        =
        (4\pi\tau)^{-d/2}\exp\!\left(-\frac{|x-y|^2}{4\tau}\right),
        \qquad
        K_\tau\to\delta
        \quad(\tau\downarrow0).
\]

\begin{remark}[Sector identity, not universal identity]
If \(P\neq I\), then the kernel is an identity only on the retained coherent sector. The limit \(\tau\downarrow0\) gives the sector kernel \(P\chi(A)^2P\), not the full-space Dirac delta unless the sector/window tends to the full identity. This distinction is essential.
\end{remark}


\section{A worked fixed-point selection example: \(S^1_R\)}

To make the selection mechanism explicit, consider the coherent chart \(X=S^1_R\), the circle of radius \(R\), with coordinate \(\theta\in[0,2\pi)\). Let
\[
        A=-R^{-2}\partial_\theta^2 .
\]
The eigenfunctions and eigenvalues are
\[
        \phi_n(\theta)=\frac{1}{\sqrt{2\pi R}}e^{in\theta},
        \qquad
        \lambda_n=\frac{n^2}{R^2},
        \qquad n\in\mathbb Z .
\]
If the fixed-point cycle supplies a proper-time scale \(\tau_0>0\), the canonical heat/proper-time kernel is not chosen by hand. It is fixed:
\[
        K_{\tau_0}^{S^1_R}(\theta,\theta')
        =
        \frac{1}{2\pi R}
        \sum_{n\in\mathbb Z}
        e^{-\tau_0 n^2/R^2}
        e^{in(\theta-\theta')}.
\]
The corresponding scalar propagator with mass \(m>0\) is
\[
        \Delta_{\tau_0}^{S^1_R}(\theta,\theta')
        =
        \frac{1}{2\pi R}
        \sum_{n\in\mathbb Z}
        \frac{e^{-\tau_0 n^2/R^2}}{n^2/R^2+m^2}
        e^{in(\theta-\theta')}.
\]
In this example the finite correction is fully determined by the triple
\[
        (R,A,\tau_0).
\]
There is no independent choice of a Gaussian regulator.

The effective spectral width is read off from
\[
        e^{-\tau_0 n^2/R^2}\sim e^{-1}
        \quad\Longleftrightarrow\quad
        |n|\sim \frac{R}{\sqrt{\tau_0}},
\]
or, in physical momentum units,
\[
        p_{\mathrm{eff}}\sim \tau_0^{-1/2}.
\]
This is the simplest explicit model of the general claim:
\[
        \boxed{
        \text{fixed-point operator data select the finite kernel.}
        }
\]


\section{MTT-corrected scalar propagator}

Let
\[
        L_m := A+m^2 = -\Delta+m^2,
        \qquad m>0.
\]
The ordinary Euclidean propagator is
\[
        \Delta_0 := L_m^{-1}.
\]
The MTT-corrected propagator is the filtered covariance
\[
        \Delta_{\tau,\chi}
        :=
        B_{\tau,\chi}L_m^{-1}B_{\tau,\chi}^{\ast}.
\]
In the full flat chart \(P=\chi=I\), this becomes
\[
        \Delta_\tau
        =
        e^{-\tau A/2}(A+m^2)^{-1}e^{-\tau A/2}
        =
        e^{-\tau A}(A+m^2)^{-1}.
\]
Hence in momentum space,
\begin{equation}
        \label{eq:filtered-propagator}
        \widehat{\Delta_\tau}(k)
        =
        \frac{e^{-\tau |k|^2}}{|k|^2+m^2}.
\end{equation}

\begin{definition}[MTT effective coherence scale]
The proper-time parameter \(\tau>0\) defines the effective coherence scale
\[
        \Lambda_{\mathrm{eff}}
        :=
        \tau^{-1/2}.
\]
Modes with \(|k|\gg \Lambda_{\mathrm{eff}}\) are exponentially suppressed.
\end{definition}

\section{Recovery of the ordinary propagator}

\begin{theorem}[Distributional recovery]
Let \(m>0\), \(d<\infty\), and define \(\Delta_\tau\) by \cref{eq:filtered-propagator}. Then
\[
        \Delta_\tau \longrightarrow \Delta_0
\]
in the sense of tempered distributions as \(\tau\downarrow0\). Equivalently, for all Schwartz functions \(f,g\in\mathcal S(\mathbb R^d)\),
\[
        \lim_{\tau\downarrow0}
        \int_{\mathbb R^d}
        \frac{e^{-\tau |k|^2}}{|k|^2+m^2}
        \widehat f(k)^\ast \widehat g(k)
        \frac{d^dk}{(2\pi)^d}
        =
        \int_{\mathbb R^d}
        \frac{\widehat f(k)^\ast \widehat g(k)}{|k|^2+m^2}
        \frac{d^dk}{(2\pi)^d}.
\]
\end{theorem}

\begin{proof}
For each \(k\), \(e^{-\tau |k|^2}\to1\). Moreover,
\[
        0\le
        \frac{e^{-\tau |k|^2}}{|k|^2+m^2}
        \le
        \frac{1}{|k|^2+m^2}.
\]
Since \(f,g\) are Schwartz, the function
\[
        \frac{|\widehat f(k)|\,|\widehat g(k)|}{|k|^2+m^2}
\]
is integrable. Dominated convergence gives the result.
\end{proof}

\section{Finite coincident value}

The ordinary coincident propagator
\[
        \Delta_0(x,x)
        =
        \int_{\mathbb R^d}\frac{d^dk}{(2\pi)^d}\frac{1}{|k|^2+m^2}
\]
diverges for \(d\ge2\). The filtered propagator has
\[
        \Delta_\tau(x,x)
        =
        I_d(m,\tau)
        :=
        \int_{\mathbb R^d}
        \frac{d^dk}{(2\pi)^d}
        \frac{e^{-\tau |k|^2}}{|k|^2+m^2}.
\]

\begin{theorem}[Finite diagonal]
For every \(d<\infty\), \(m>0\), and \(\tau>0\),
\[
        I_d(m,\tau)<\infty .
\]
Moreover,
\[
        I_d(m,\tau)
        =
        \frac{1}{(4\pi)^{d/2}}
        \int_{\tau}^{\infty}
        u^{-d/2}e^{-m^2u}\,du .
\]
\end{theorem}

\begin{proof}
Use
\[
        \frac{1}{|k|^2+m^2}
        =
        \int_0^\infty e^{-s(|k|^2+m^2)}\,ds.
\]
Then
\[
        I_d(m,\tau)
        =
        \int_0^\infty e^{-m^2s}
        \left[
        \int_{\mathbb R^d}
        \frac{d^dk}{(2\pi)^d}
        e^{-(s+\tau)|k|^2}
        \right]ds.
\]
The Gaussian integral gives
\[
        \int_{\mathbb R^d}
        \frac{d^dk}{(2\pi)^d}
        e^{-(s+\tau)|k|^2}
        =
        (4\pi(s+\tau))^{-d/2}.
\]
With \(u=s+\tau\),
\[
        I_d(m,\tau)
        =
        \frac{1}{(4\pi)^{d/2}}
        \int_{\tau}^{\infty}
        u^{-d/2}e^{-m^2(u-\tau)}\,du .
\]
If the filter is placed as \(e^{-\tau A/2}\) on each side, the propagator carries \(e^{-\tau |k|^2}\), giving the harmless factor \(e^{m^2\tau}\) depending on the chosen proper-time normalization. Equivalently one may define the shifted proper-time representation
\[
        \frac{e^{-\tau |k|^2}}{|k|^2+m^2}
        =
        e^{m^2\tau}
        \int_{\tau}^{\infty}
        e^{-u(|k|^2+m^2)}\,du.
\]
Thus
\[
        I_d(m,\tau)
        =
        \frac{e^{m^2\tau}}{(4\pi)^{d/2}}
        \int_{\tau}^{\infty}
        u^{-d/2}e^{-m^2u}\,du .
\]
Both forms are finite for \(\tau>0\). The displayed normalization depends on whether the mass term is included in the heat filter. In either convention the lower proper-time gap \(\tau>0\) removes the \(u=0\) singularity.
\end{proof}

\begin{remark}[Normalization convention]
Two closely related filters occur in the literature:
\[
        e^{-\tau A}(A+m^2)^{-1},
        \qquad
        e^{-\tau(A+m^2)}(A+m^2)^{-1}.
\]
They differ by the finite factor \(e^{-\tau m^2}\). All UV statements are unchanged. When positivity/proper-time support is emphasized, the second convention is often cleaner.
\end{remark}

\section{Leading low-energy correction}

For \(\tau |k|^2\ll1\),
\[
        e^{-\tau |k|^2}
        =
        1-\tau |k|^2+\frac{\tau^2|k|^4}{2}+O(\tau^3|k|^6).
\]
Therefore
\[
        \widehat{\Delta_\tau}(k)
        =
        \frac{1}{|k|^2+m^2}
        -
        \tau \frac{|k|^2}{|k|^2+m^2}
        +
        O\!\left(\tau^2\frac{|k|^4}{|k|^2+m^2}\right).
\]

\begin{proposition}[Low-momentum correction bound]
For \(|k|^2\le R^2\),
\[
        \left|
        \frac{e^{-\tau |k|^2}-1}{|k|^2+m^2}
        \right|
        \le
        \frac{\tau |k|^2}{|k|^2+m^2}
        \le
        \tau .
\]
Thus the pointwise low-energy correction is \(O(\tau |k|^2)\) relative to the local numerator and bounded by \(O(\tau)\) in absolute propagator units on compact momentum sets.
\end{proposition}

\begin{proof}
Use \(0\le 1-e^{-a}\le a\) for \(a\ge0\).
\end{proof}

\section{One-loop examples}

\subsection{Tadpole}

The ordinary one-loop tadpole integral is
\[
        T_0
        =
        \int_{\mathbb R^d}
        \frac{d^dk}{(2\pi)^d}
        \frac{1}{|k|^2+m^2},
\]
which is divergent for \(d\ge2\). The MTT-filtered tadpole is
\[
        T_\tau
        =
        \int_{\mathbb R^d}
        \frac{d^dk}{(2\pi)^d}
        \frac{e^{-\tau |k|^2}}{|k|^2+m^2},
\]
which is finite by the finite-diagonal theorem.

\subsection{Bubble}

For external momentum \(p\),
\[
        B_\tau(p)
        =
        \int_{\mathbb R^d}
        \frac{d^dk}{(2\pi)^d}
        \frac{e^{-\tau |k|^2}e^{-\tau |k+p|^2}}
        {(|k|^2+m^2)(|k+p|^2+m^2)} .
\]

\begin{proposition}[Bubble finiteness]
For every \(p\in\mathbb R^d\), \(m>0\), and \(\tau>0\),
\[
        B_\tau(p)<\infty .
\]
\end{proposition}

\begin{proof}
For large \(|k|\), \(|k+p|^2\ge \frac12|k|^2-|p|^2\). Hence
\[
        e^{-\tau |k|^2}e^{-\tau |k+p|^2}
        \le
        C_p e^{-c_\tau |k|^2}
\]
for some \(C_p,c_\tau>0\). The denominator is positive and polynomial, so the integrand is dominated by an integrable Gaussian.
\end{proof}

\section{All-graph Euclidean UV domination}

We now state a controlled perturbative result for scalar Euclidean diagrams. The point is not to construct a full QFT, but to show that canonical MTT filtering gives absolute UV convergence once every internal line carries the fixed-point Gaussian factor.

Consider a connected graph \(\Gamma\) with \(L\) loop momenta \(\ell=(\ell_1,\dots,\ell_L)\in(\mathbb R^d)^L\), internal lines \(e\in E(\Gamma)\), and fixed external momenta. Each internal line momentum has the affine form
\[
        q_e(\ell)=M_e\ell+r_e,
\]
where \(M_e\) is a linear map and \(r_e\) depends on external momenta. The filtered scalar integrand is bounded by
\[
        \prod_{e\in E(\Gamma)}
        \frac{e^{-\tau |q_e(\ell)|^2}}
        {|q_e(\ell)|^2+m^2}
\]
times a polynomial numerator \(P(\ell)\), if derivative interactions are present.

\begin{lemma}[Loop-momentum quadratic control]
For a connected graph and a choice of independent loop momenta, there exist constants \(c_\Gamma>0\) and \(C_\Gamma\ge0\), depending on the graph and external momenta, such that
\[
        \sum_{e\in E(\Gamma)} |q_e(\ell)|^2
        \ge
        c_\Gamma |\ell|^2-C_\Gamma .
\]
\end{lemma}

\begin{proof}
If the inequality failed, there would be a sequence \(|\ell_n|\to\infty\) with all line momenta \(q_e(\ell_n)\) sublinear in \(|\ell_n|\). After normalizing, a nonzero loop direction would lie in the kernel of every line-momentum map \(M_e\). But independent loop momenta are chosen precisely so that each nonzero loop flow crosses at least one internal line. Hence the common kernel is trivial. Finite-dimensional norm equivalence then gives the stated coercivity, with the affine external shifts absorbed into \(C_\Gamma\).
\end{proof}

\begin{theorem}[Euclidean all-graph UV finiteness under canonical line filtering]
Fix \(m>0\), \(\tau>0\), and a connected scalar Euclidean Feynman graph \(\Gamma\). Suppose every internal line carries the canonical Gaussian factor \(e^{-\tau |q_e|^2}\). Then the loop integral
\[
        \int_{(\mathbb R^d)^L}
        |P(\ell)|
        \prod_{e\in E(\Gamma)}
        \frac{e^{-\tau |q_e(\ell)|^2}}
        {|q_e(\ell)|^2+m^2}
        \,d^{dL}\ell
\]
is finite for every polynomial numerator \(P\).
\end{theorem}

\begin{proof}
By the quadratic-control lemma,
\[
        \prod_e e^{-\tau |q_e(\ell)|^2}
        =
        e^{-\tau\sum_e |q_e(\ell)|^2}
        \le
        e^{\tau C_\Gamma}e^{-\tau c_\Gamma|\ell|^2}.
\]
The denominator is bounded above by dropping it, since \((|q_e|^2+m^2)^{-1}\le m^{-2}\). Thus the absolute integrand is bounded by
\[
        C |P(\ell)|e^{-\tau c_\Gamma|\ell|^2},
\]
which is integrable on finite-dimensional loop-momentum space.
\end{proof}

\begin{remark}[Gauge and Lorentzian caveat]
This theorem is Euclidean and scalar. Gauge theories require compatible filtering of fields, ghosts, vertices, and Ward/BRST identities. Lorentzian unitarity requires a separate analytic-continuation or positive spectral representation argument. The present theorem proves UV domination, not by itself full physical consistency.
\end{remark}

\section{Proper-time representation and positivity}

For the mass-including convention,
\[
        \Delta_\tau^{(m)}(k)
        =
        \frac{e^{-\tau(|k|^2+m^2)}}{|k|^2+m^2},
\]
one has the exact proper-time representation
\[
        \Delta_\tau^{(m)}(k)
        =
        \int_{\tau}^{\infty} e^{-s(|k|^2+m^2)}\,ds .
\]
This is a positive proper-time integral with a lower support gap \(s\ge\tau\). The ordinary propagator is recovered by lowering the gap to zero:
\[
        \frac{1}{|k|^2+m^2}
        =
        \int_{0}^{\infty} e^{-s(|k|^2+m^2)}\,ds.
\]

This representation gives the MTT reading:
\[
        \boxed{\text{UV softening}=\text{absence of arbitrarily short proper-time support}.}
\]
It matches the fixed-point interpretation of \(\tau\) as finite coherent support rather than an arbitrary regulator.

\section{Regulator versus MTT kernel}

Mathematically, the factor \(e^{-\tau |k|^2}\) resembles a UV regulator. The MTT claim is stronger and narrower:
\[
        \boxed{
        \text{the filter is admissible only when derived from fixed-point operator data.}
        }
\]
The direction of explanation is therefore reversed relative to ordinary regularization. In a regulator calculation one first has a divergent integral and then chooses a cutoff. In the MTT execution reading one first has fixed-point data \((A,P,\chi,\tau)\), obtains the coherent filter, and only then computes amplitudes. The Gaussian factor is a consequence of coherent stabilization, not a post hoc cure for a divergence.

An arbitrary smearing function may break locality, positivity, gauge symmetry, or unitarity. A physically admissible MTT kernel must be selected by the coherent projector, spectral window, proper-time support, and compatibility conditions of the relevant sector.

\begin{center}
\begin{tabular}{p{0.34\linewidth}p{0.52\linewidth}}
\toprule
Object & Role \\
\midrule
\(A\) & fixed-point linearized damping/geometric operator \\
\(P\) & coherent-sector projector \\
\(\chi(A)\) & admissible spectral window \\
\(\tau\) & proper-time/coherence support scale \\
\(B_{\tau,\chi}\) & canonical coherent filter \\
\(\Delta_{\tau,\chi}\) & MTT-corrected covariance/propagator \\
\(\Lambda_{\mathrm{eff}}\sim\tau^{-1/2}\) & effective UV/coherence scale \\
\bottomrule
\end{tabular}
\end{center}


\section{Numerical anchoring from existing data}

The model does not yet derive \(\tau\) numerically from the full MTT carrier. Nevertheless, once the propagator correction is fixed, existing data translate directly into bounds or target scales.

For low momenta,
\[
        e^{-\tau k^2}=1-\tau k^2+O(\tau^2 k^4).
\]
Thus an experiment that sees no deviation from the local propagator up to characteristic momentum scale \(E\), with relative tolerance \(\varepsilon_{\mathrm{exp}}\), gives the parametric bound
\[
        \tau E^2 \lesssim \varepsilon_{\mathrm{exp}},
        \qquad
        \Lambda_{\mathrm{eff}}=\tau^{-1/2}
        \gtrsim \frac{E}{\sqrt{\varepsilon_{\mathrm{exp}}}}.
\]
This is deliberately written as a bound rather than a prediction, because the sector value of \(\tau\) is not fixed here.

As an illustrative collider-scale anchoring, LHC Run 3 operates at a proton--proton collision energy of \(13.6\,\mathrm{TeV}\) \cite{CERNRun3}. Requiring the coherence scale to lie above this scale gives the conservative estimate
\[
        \Lambda_{\mathrm{eff}}\gtrsim 13.6\,\mathrm{TeV},
        \qquad
        \sqrt{\tau}\lesssim 1.45\times 10^{-20}\,\mathrm{m}.
\]
A stronger illustrative benchmark comes from contact-interaction/compositeness searches, where global summaries commonly quote lower scales in the multi-TeV to tens-of-TeV range \cite{PDGCompositeness2025}. Taking
\[
        \Lambda_{\mathrm{eff}}\gtrsim 20\,\mathrm{TeV}
        \quad\text{or}\quad
        50\,\mathrm{TeV}
\]
corresponds to
\[
        \sqrt{\tau}\lesssim 9.87\times10^{-21}\,\mathrm{m},
        \qquad
        \sqrt{\tau}\lesssim 3.95\times10^{-21}\,\mathrm{m},
\]
respectively. These numbers should not be read as established MTT predictions. They show how existing high-energy data would bound the coherence length once the filtered-propagator correction is identified with a physical sector.

For comparison, the 2022 CODATA Planck length is
\[
        \ell_{\mathrm P}=1.616255(18)\times10^{-35}\,\mathrm{m}
\]
\cite{NISTPlanck}. A Planck-scale coherence length would correspond to
\[
        \Lambda_{\mathrm{eff}}\sim 1.22\times10^{19}\,\mathrm{GeV},
\]
far beyond present collider reach. The present framework allows either possibility: \(\tau\) may be experimentally bounded, or it may be fixed by deeper MTT carrier data at a much smaller length scale.

\begin{center}
\begin{tabular}{p{0.28\linewidth}p{0.28\linewidth}p{0.30\linewidth}}
\toprule
Benchmark & \(\Lambda_{\mathrm{eff}}\) & \(\sqrt{\tau}\sim \Lambda_{\mathrm{eff}}^{-1}\) \\
\midrule
LHC Run 3 collision scale & \(13.6\,\mathrm{TeV}\) & \(1.45\times10^{-20}\,\mathrm{m}\) \\
Illustrative contact scale & \(20\,\mathrm{TeV}\) & \(9.87\times10^{-21}\,\mathrm{m}\) \\
Illustrative contact scale & \(50\,\mathrm{TeV}\) & \(3.95\times10^{-21}\,\mathrm{m}\) \\
Planck length benchmark & \(1.22\times10^{19}\,\mathrm{GeV}\) & \(1.62\times10^{-35}\,\mathrm{m}\) \\
\bottomrule
\end{tabular}
\end{center}

The important point is conceptual: the filtered propagator converts any empirical lower bound on unresolved short-distance structure into a bound on the MTT coherence scale,
\[
        \sqrt{\tau}\lesssim \Lambda_{\mathrm{eff}}^{-1}.
\]


\section{What becomes predictive}

The paper proves that, once \(A,P,\chi,\tau\) are fixed by an MTT sector, the propagator correction is fixed. The remaining execution tasks are now concrete:
\[
        (A,P,\chi,\tau)
        \Longrightarrow
        \Delta_{\tau,\chi}
        \Longrightarrow
        \text{observable correction}.
\]
Examples include:
\begin{enumerate}
\item finite coincident correlators;
\item modified short-distance potentials;
\item exponentially softened loop amplitudes;
\item low-energy corrections of order \(\tau |k|^2\);
\item effective scale constraints \(\Lambda_{\mathrm{eff}}\sim\tau^{-1/2}\).
\end{enumerate}

The crucial change from the earlier dictionary papers is that the kernel is no longer freely chosen. The corrected propagator is dictated by the canonical fixed-point filter.

\section{Conclusion}

The delta-projection program required a transition from interpretation to execution. The canonical coherent-kernel paper supplied the missing kernel-selection theorem. This paper applied that theorem to propagators.

In a flat coherent chart, MTT fixed-point data select the corrected propagator
\[
        \widehat{\Delta_\tau}(k)
        =
        \frac{e^{-\tau |k|^2}}{|k|^2+m^2}.
\]
This propagator recovers the ordinary local propagator as \(\tau\downarrow0\), has finite coincident value for every finite dimension when \(\tau>0\), and gives all-graph Euclidean UV domination for scalar perturbation theory when every internal line carries the canonical filter.

The result does not yet derive the numerical value of \(\tau\) from the full MTT carrier. It does something more basic and necessary: it shows that once an MTT fixed-point sector supplies \(\tau\) and the operator data, the finite correction is specific. Existing high-energy data can then be translated into bounds on \(\sqrt{\tau}\), as illustrated above.

\[
        \boxed{
        \text{MTT does not merely soften deltas; it selects the propagator correction.}
        }
\]

\begin{thebibliography}{9}

\bibitem{ReedSimon}
M. Reed and B. Simon,
\emph{Methods of Modern Mathematical Physics}.
Academic Press.

\bibitem{SimonFunctional}
B. Simon,
\emph{Functional Integration and Quantum Physics}.
AMS Chelsea.

\bibitem{GlimmJaffe}
J. Glimm and A. Jaffe,
\emph{Quantum Physics: A Functional Integral Point of View}.
Springer.

\bibitem{Rivasseau}
V. Rivasseau,
\emph{From Perturbative to Constructive Renormalization}.
Princeton University Press.


\bibitem{CERNRun3}
CERN,
``The third run of the Large Hadron Collider has successfully started,''
CERN News, 5 July 2022.

\bibitem{PDGCompositeness2025}
Particle Data Group,
``Searches for Quark and Lepton Compositeness,''
Review of Particle Physics, revised August 2025.

\bibitem{NISTPlanck}
NIST CODATA,
``CODATA Value: Planck length,''
2022 CODATA recommended values.


\bibitem{Kleinert}
H. Kleinert,
\emph{Path Integrals in Quantum Mechanics, Statistics, Polymer Physics, and Financial Markets}.
World Scientific.

\end{thebibliography}

\end{document}
